\begin{textsample}{POS Dim 3 – am – Score 31.00 – am/am\_inf\_41.txt}  \label{ex:f3_pos_035}
20.
Estrutura organizacional das \textbf{creches} e pré-escolas O espaço físico das unidades de ensino deve priorizar um atendimento que garanta o direito estabelecido a todas as \textbf{crianças} que nele transitam, inclusive aquelas com deficiência, transtornos globais do desenvolvimento e altas habilidades/superdotação, que tenham assegurados o direito ao \textbf{cuidado} e à educação, à saúde, proteção, descanso, interação, \textbf{conforto}, \textbf{higiene} e aconchego.
No atendimento para as fases \textbf{creche} e \textbf{pré-escola}, devem ser observadas a concepção da edificação, vislumbrando a funcionalidade e acessibilidade esperada para esses ambientes, respeitando os critérios definidos pela Associação Brasileira de Normas Técnicas – ABNT, NBR 9050/2004 e regulamentados pela lei n.º 10.098/2000, abrangendo : organização espacial e dimensionamento dos conjuntos funcionais, acessibilidade, percursos, segurança e adequação \textbf{mobiliária}.
Conforme as orientações descritas nos Parâmetros Básicos de Infraestrutura para \textbf{Instituição} de Educação \textbf{Infantil} ( 2006b ), os aspectos estéticos dos espaços dizem respeito à imagem e à aparência do ambiente, proporcionando momentos diferenciados e de bem-estar.
Inclui -se as variedades e diversidades de \textbf{texturas} e cores, suas formas, proporções e símbolos, valorizando os elementos visuais da edificação, que podem nortear os trabalhos e atividades, visando o despertar dos sentidos, a \textbf{curiosidade} e a capacidade de descoberta do \textbf{bebê}, da \textbf{criança} bem \textbf{pequena} e da \textbf{criança} \textbf{pequena}.
Durante a seleção destes materiais devem ser considerados os seguintes aspectos : - Qualidade técnica : é considerada por sua eficiência, função, facilidade de \textbf{manutenção} e limpeza ; - Qualidade ergonômica : direciona -se a facilidade de \textbf{manuseio} do produto, adaptação antropométrica, clareza de informações, compatibilidade de movimentos, \textbf{conforto} e segurança ; - Qualidade estética : engloba toda combinação das formas, cores, uso de materiais e \textbf{texturas}, proporcionando uma \textbf{agradável} leveza em sua aparência.
Os espaços internos das salas de referência devem propiciar às \textbf{crianças} uma ampla \textbf{movimentação}.
Estes espaços devem ser organizados e livres de itens que possam oferecer riscos físicos às \textbf{crianças} ( como tomadas baixas, que deverão receber tampa protetora ), \textbf{contar} com uma boa iluminação, ser ventilados, estar em uma temperatura ambiente confortável, conservados e limpos com frequência.
A limpeza dos espaços deve ser sistemática, seguindo orientações pertinentes à vigilância em saúde nos espaços coletivos de atendimento de \textbf{bebês}, \textbf{crianças} bem \textbf{pequenas} e \textbf{crianças} \textbf{pequenas}.
O vestuário e o calçado das pessoas responsáveis por tal ação devem atender às normas de segurança e os produtos usados devem ser adequados, de boa qualidade e jamais deixados ao alcance das \textbf{crianças}.
O espaço físico também educa.
Como nossos espaços estão organizados para \textbf{acolher} as \textbf{crianças} e educar seu olhar?
O que existe em nossas \textbf{paredes} que contribui para a sensação de pertencimento da \textbf{criança} a esse espaço?
Para isso, precisamos observar como os \textbf{adultos} entendem e organizam os espaços.
Algumas questões para considerar - As janelas devem estar na altura que permitam às \textbf{crianças} a visualização do espaço externo, o material de ambientação da sala de referência ( cartazes, lista com nome e aniversários, \textbf{calendário}, quadro do tempo, fotografias, desenhos, \textbf{pinturas} ( dos artistas e das \textbf{crianças}! ), entre outros materiais com os quais elas possam interagir ( não limitar a oferta de imagens a \textbf{personagens} de desenhos e filmes \textbf{infantis}, mas desafiar o olhar dos \textbf{adultos} e \textbf{crianças} ), devem estar na altura dos seus olhos, ao alcance de suas \textbf{mãos}, e não no alto para que elas não toquem.
As \textbf{paredes} devem servir como espaço de exposição das produções das próprias \textbf{crianças}.
O momento da ambientação da sala deve envolver tanto o professor quanto a \textbf{criança}, considerando as experiências desenvolvidas a partir do \textbf{cuidar} e educar, das interações e \textbf{brincadeiras}.
A construção de materiais para compor as \textbf{paredes} da sala de referência e dos espaços externos deve sempre \textbf{contar} com a participação das \textbf{crianças}, seja nas \textbf{sugestões}, seja na elaboração. - É muito importante que as salas de referências sejam um espaço que provoca os sentidos.
Para isso, os materiais precisam ser organizados de forma acessível e que desafie aquilo que a \textbf{criança} constrói a partir de suas hipóteses sobre o mundo.
Os \textbf{livros} de literatura, \textbf{brinquedos} estruturados e não estruturados devem compor o ambiente, além de materiais para o exercício da criação : desenho, \textbf{pintura}, modelagem e construção, o que não significa que essas atividades só ocorrem se houver materiais industrializados.[6 ] - Além das salas de referência, é necessário que existam outros espaços que ampliem a exploração de recursos pedagógicos, como \textbf{bibliotecas}, brinquedotecas e laboratórios.
Quando houver, a organização desses ambientes deve ser coletiva, considerando a opinião e a construção da \textbf{criança} nesse processo, devendo ainda gerar discussões e ideias reflexivas ao grupo, favorecendo às \textbf{crianças} um ambiente institucional mais atraente, acolhedor, instigante e propulsor de novos mecanismos de aprendizagens entre criança/criança, criança/adultos e criança/ambiente.
A criação desses espaços deve ter como foco a funcionalidade e estar em conformidade com a proposta pedagógica e o planejamento de cada unidade de ensino. [ ^6 ] : Conferir o texto As múltiplas linguagens, no capítulo 7, p. 41.

% matched lemmas: acolher, adulto, agradável, bebê, biblioteca, brincadeira, brinquedo, calendário, conforto, contar, creche, criança, cuidado, cuidar, curiosidade, higiene, infantil, instituição, livro, manuseio, manutenção, mobiliário, movimentação, mão, parede, pequeno, personagem, pintura, pré-escola, sugestão, textura
\end{textsample}
